\chapter{Basic sets and equivariant bijections}
\label{chap:library:assembly}

An integral basic set identifies two free abelian groups, one indexed
by ordinary characters and the other by Brauer characters. If a group
acts compatibly with decomposition, this identification is equivariant.
For a hypoelementary group, Conlon's theorem and injectivity of the mark homomorphism
give a bijection of the two sets of characters. Further results
in this chapter extend equivariant maps from stabilisers, cancel
invariant subsets and construct fixed representatives in product orbits.

Write $A$ for an acting group. An $A$-equivariant map satisfies
$f(ax)=af(x)$, and $A_x$ denotes the stabiliser of $x$. If a map $\pi:X\to I$
is given, its fibre over $i$ is the subset $X_i=\{x:\pi(x)=i\}$.

\section{Naturality and an integral basic set}

Let $G$ be finite and let $\ell$ be prime. Let $K$ and $k$ be fields
of characteristic zero and $\ell$, respectively. Consider the Grothendieck groups
$K_0(KG)$ and $K_0(kG)$ of Chapter~\ref{chap:library:kzero}, and a specified
decomposition homomorphism $d:K_0(KG)\to K_0(kG)$. Let $X$ and $Y$
be finite sets. For a set $X$, let
$\Z X$ denote the free abelian group with basis $(e_x)_{x\in X}$.
Suppose $X$ and $Y$ index distinct simple modules over $KG$ and $kG$,
respectively. Sending a basis vector to the corresponding module class
gives injective maps $j_X:\Z X\to K_0(KG)$ and
$j_Y:\Z Y\to K_0(kG)$, by the basis of simple classes in
Theorem~\ref{lib:kzero:basis-theorem}.

\begin{definition}
The family indexed by $X$ is an \emph{integral basic set} for the family
indexed by $Y$ if the restriction of $d$ is an isomorphism
$j_X(\Z X)\to j_Y(\Z Y)$. Equivalently, there is an isomorphism
$D:\Z X\to\Z Y$ such that $j_YD=dj_X$.
If $A$ acts on $X$ and $Y$,
it acts on the free groups by $a e_x=e_{ax}$ and $a e_y=e_{ay}$.
Compatible actions on the two Grothendieck groups are actions for which
$j_X$ and $j_Y$ commute with $A$.
\end{definition}

\begin{proposition}[Equivariance of the restricted decomposition map]
\label{lib:actions:basic-set}
Suppose these actions are compatible and $d$ commutes with $A$. Then $D$
is an isomorphism of integral permutation lattices. If
$D(e_x)=\sum_y d_{y,x}e_y$, its coefficients satisfy
$d_{ay,ax}=d_{y,x}$ for every $a\in A$.
\end{proposition}
\begin{proof}
For $v\in\Z X$, use the defining square and naturality to obtain
\[
 j_Y D(av)=d j_X(av)=d(a j_Xv)
          =a d j_Xv=a j_Y Dv=j_Y(a Dv).
\]
Injectivity of $j_Y$ gives $D(av)=aD(v)$. Taking coefficients on basis
vectors proves the stated matrix identity. Conversely, that identity on
basis vectors implies equivariance by linearity.
\end{proof}

The next section
gives conditions under which the two indexing sets are also
equivariantly bijective. For automorphism twisting, the construction in
Chapter~\ref{chap:library:reduction} gives naturality under its stated
compatibility hypothesis. For tensoring with linear characters, the ordinary and modular actions
must satisfy the corresponding character and root compatibility laws.
\begin{implementation}
The permutation representations are Mathlib's
\lean{Representation.ofMulAction}, with coefficients in $\Z$.
\module{ModularRep.DecompositionBasicSetBridge} defines
\lean{RestrictedIntegralBasicSet} and proves
\lean{matrixEquivariant_of_kZero_naturality}. Its
\lean{matrixEquivariant_of_stableReduction} obtains naturality from the
constructed decomposition map. The equivalence between the matrix and
linear map formulations is
in \module{ModularRep.IntegralBasicSetBridge}.
\end{implementation}

\begin{proposition}[Restriction to an invariant block]
\label{lib:actions:block-restriction}
Suppose $D$ is equivariant and $\pi_X:X\to I$, $\pi_Y:Y\to I$ are
$A$-invariant block assignments. If
\[
 D(\Z X_i)\subseteq\Z Y_i,
 \qquad D^{-1}(\Z Y_i)\subseteq\Z X_i
 \qquad(i\in I),
\]
then $D$ restricts to an isomorphism of $\Z A$-lattices
$\Z X_i\to\Z Y_i$ for every $i\in I$.
\end{proposition}
\begin{proof}
The support conditions restrict $D$ and its inverse to the two submodules
spanned by the indicated fibres. The restrictions are inverse because the
original maps are inverse. The action preserves each submodule, and the
restricted maps still commute with it. Identifying each supported submodule
with the free group on its fibre gives the assertion.
\end{proof}
\begin{implementation}
\module{ModularRep.BlockFibreRestriction}, especially
\lean{restrictBlockLinearEquiv}, \lean{restrictBlockLinearEquiv_equivariant}
and \lean{restrictRestrictedIntegralBasicSet}. Both support conditions
are hypotheses.
\end{implementation}

\section{From permutation lattices to finite sets}

Let $p$ be a prime, and write $\Z_p$ for the ring of $p$-adic
integers. The prime $p$ may differ from the characteristic $\ell$
used for modular characters. A finite group $A$ is
\emph{$p$-hypoelementary} if it has a normal $p$-subgroup with cyclic
quotient of order prime to $p$.
For an $A$-set $X$ and a subgroup $U\leq A$, the \emph{mark} at $U$
is $|X^U|$, where
\[
 X^U=\{x\in X:ux=x\text{ for every }u\in U\}.
\]

\begin{theorem}[Conlon]
Let $A$ be a finite group and let $X,Y$ be finite $A$-sets.
If $\Z_p X$ and $\Z_p Y$ are isomorphic as $\Z_p A$-lattices, then
\[
 |X^U|=|Y^U|
\]
for every $p$-hypoelementary subgroup $U$ of $A$.
\end{theorem}

Conlon's fixed coset calculation gives this implication
\cite[Section~4, especially (4.7)]{Conlon1968}.
Bouc gives a modern formulation, including the converse, in
\cite[Theorem~3.5.5]{Bouc2000}.
In that account the coefficient ring is a complete commutative
Noetherian local ring with residue characteristic $p$, and lattices
are finitely generated and free over that ring. The choice $\Z_p$
satisfies these hypotheses.

\begin{theorem}[Burnside's determination by marks]
Let $A$ be a finite group. Two finite $A$-sets $X$ and $Y$ are
isomorphic if and only if $|X^U|=|Y^U|$ for every subgroup $U\leq A$.
\end{theorem}
\begin{proof}
An equivariant bijection identifies the fixed points of every subgroup.
For the converse, decompose each set into transitive orbits $A/H$,
with $H$ taken from a set of representatives of subgroup conjugacy
classes. The matrix
\[
 m(U,H)=|(A/H)^U|
\]
is triangular when those representatives are ordered by subgroup
order. Its diagonal entries are $|N_A(H):H|$, which are nonzero.
Equality of marks therefore gives equality of the multiplicities of
all transitive orbits.
\end{proof}
For this statement and proof, see, for example,
\cite[Theorem~2.3.2]{Bouc2000}.
The equivalent formulation is injectivity of the mark homomorphism
from the Burnside ring.

\begin{corollary}[Permutation lattices]
\label{lib:actions:conlon}
Let $A$ be a finite $p$-hypoelementary group and let $X,Y$ be finite
$A$-sets. If $\Z_p X\simeq\Z_p Y$ as $\Z_p A$-lattices, then
$X$ and $Y$ are equivariantly bijective. In particular, the same
conclusion holds if $\Z X\simeq\Z Y$ as $\Z A$-lattices.
\end{corollary}
\begin{proof}
Conlon's theorem gives equality of the marks at all
$p$-hypoelementary subgroups. Every subgroup of $A$ is
$p$-hypoelementary by Proposition~\ref{lib:groups:hypoelementary}.
Thus all marks agree, and the preceding theorem gives the bijection.
For the integral version, extend the isomorphism and its inverse
from $\Z$ to $\Z_p$.
\end{proof}
\begin{implementation}
\module{ModularRep.PaperProofs.ConlonBasicSet} constructs scalar
extension and the mark map on isomorphism classes of finite $A$-sets.
The scalar extension uses Mathlib's equivalences for tensor products
with finitely supported functions. The intertwining identities are
proved in this module.
The two published theorems are separate hypotheses:
\lean{PadicConlonMarkDetection} and
\lean{PublishedBurnsideMarkInjectivity}. Both range over all finite
$A$-sets. Under these hypotheses,
\lean{corollary_2_4_conlonMark} proves the implication from
isomorphic $p$-adic permutation lattices to an equivariant bijection.
The proofs of the published theorems are not part of that declaration.
\end{implementation}

If $A$ is $p$-hypoelementary, Proposition~\ref{lib:actions:basic-set}
and Corollary~\ref{lib:actions:conlon} give an equivariant bijection
between the ordinary basic set and the modular characters it indexes.
Under the same hypothesis, Proposition~\ref{lib:actions:block-restriction}
gives the corresponding bijection on an invariant block.
\begin{implementation}
The corresponding declarations are
\lean{equivariantSetEquiv_of_kZero_naturality} and
\lean{equivariantSetEquiv_restrictBlock_of_kZero_naturality}
in the modules for basic sets and block restriction.
They assume the specified integral basic set, actions, naturality,
hypoelementary group and published mark theorems.
They prove existence of a bijection without prescribing its pairs.
No unitriangularity hypothesis is used.
\end{implementation}

\section{Extending a bijection over an orbit}

Let $\pi_X:X\to I$ and
$\pi_Y:Y\to I$ be equivariant maps, and assume $I$ is transitive.

\begin{proposition}[Extension from one fibre]
\label{lib:actions:fibre-transport}
Fix $i_0\in I$ and an $A_{i_0}$-equivariant bijection
$\beta_0:X_{i_0}\to Y_{i_0}$. There is a unique $A$-equivariant bijection
$\beta:X\to Y$ that preserves the maps to $I$ and restricts to $\beta_0$ on
$X_{i_0}$.
\end{proposition}
\begin{proof}
For each $i\in I$, choose $t_i\in A$ with $t_i i_0=i$, and define
\[
 \beta(x)=t_i \beta_0(t_i^{-1}x)\qquad(x\in X_i).
\]
The expression has image in $Y_i$. Replacing $\beta_0$ by its inverse gives
an inverse map on every fibre and hence on the total sets. At $i=i_0$,
the chosen $t_{i_0}$ belongs to $A_{i_0}$, so equivariance of $\beta_0$
shows that the restriction is exactly $\beta_0$.

Two choices of $t_i$ differ by an element of $A_{i_0}$. Equivariance
of $\beta_0$ under that stabiliser makes the formulas identical. For
$a\in A$, the element $t_{ai}^{-1}at_i$ also fixes $i_0$. Applying
$A_{i_0}$-equivariance to this element proves $\beta(ax)=a\beta(x)$.
Finally, an equivariant extension of $\beta_0$ must send
$t_i(t_i^{-1}x)$ to $t_i \beta_0(t_i^{-1}x)$, which proves uniqueness.
\end{proof}
\begin{implementation}
\module{Formalisation.FibreTransport}, with \lean{Formalisation.transportedEquiv},
\lean{Formalisation.transportedEquiv_equivariant},
\lean{Formalisation.transportedEquiv_independent_of_transporter} and
\lean{Formalisation.transportedEquiv_preserves_base}.
The equivalence of the total sets uses Mathlib's \lean{Equiv.ofFiberEquiv}.
The transporter formula, its equivariance and its independence of
the choices are proved in this module.
The uniqueness assertion follows from the formula for $\beta$, as in the
proof above. It is not a separate theorem in this module.
\end{implementation}

For an arbitrary $A$-set $I$, choose one representative $i$ of each
orbit and an $A_i$-equivariant bijection $X_i\to Y_i$. Applying the
proposition on each orbit gives an $A$-equivariant bijection $X\to Y$
preserving the maps to $I$. The module
\module{Formalisation.IndexedAssembly} combines families satisfying
$\beta_{ai}(ax)=a\beta_i(x)$. The transporter formula above proves precisely
this compatibility from the bijections on the representative fibres.

The same formula preserves any specified $A$-invariant relation
$\mathcal R\subseteq X\times Y$, provided the initial bijections
satisfy it. More explicitly, assume
\begin{align*}
 (x,y)\in\mathcal R&\ \Longrightarrow\ (ax,ay)\in\mathcal R
   &&(a\in A),\\
 (x,\beta_i(x))&\in\mathcal R &&(x\in X_i)
\end{align*}
with the second condition imposed on each representative fibre.
Every $x\in X$ has the form $au$ for
some $u$ in one of those fibres, and
$(x,\beta(x))=(au,a\beta_i(u))\in\mathcal R$. Thus the local relation and
its invariance are separate hypotheses in an application.

\begin{example}[A block correspondence on an orbit]
Suppose $I$ is an orbit of blocks, $X_i$ is its Brauer character set and
$Y_i$ its set of weight classes. An equivariant bijection for one block,
under that block's stabiliser, extends to the whole orbit by
Proposition~\ref{lib:actions:fibre-transport}. It still sends each $X_i$
into $Y_i$. A modular character triple relation is preserved by the
extension provided it holds for each pair on the chosen fibre and is
invariant under the automorphisms being used.
\end{example}

\section{Involutions and cancellation}

\begin{proposition}[Classification of finite sets with an involution]
\label{lib:actions:involutions}
Let $\sigma$ and $\tau$ be involutions on finite sets $X$ and $Y$.
There is a bijection $\beta:X\to Y$ with $\beta\sigma=\tau \beta$ if and only if
$|X|=|Y|$ and $|X^\sigma|=|Y^\tau|$.
\end{proposition}
\begin{proof}
Each orbit of an involution has size one or two. If $f$ is the number
of fixed points and $t$ the number of orbits of size two, then
$|X|=f+2t$. Thus the two stated equalities imply that the numbers of
orbits of each size agree. Match the singleton orbits and the orbits of
size two. A choice of one point in each paired orbit of size two determines
the bijection on its other point by equivariance. Combining these maps gives
$\beta$. Conversely, any such $\beta$ preserves total size and carries the fixed
points bijectively to the fixed points.
\end{proof}

Equivalently, the two counts determine the cycle type of an involution
and hence its conjugacy class in the symmetric group.
\begin{implementation}
\module{Formalisation.C2Cancellation}, including
\lean{card_eq_fixed_add_twice_cycleCount} and
\lean{exists_equivariantEquiv_of_card_eq_of_fixed_card_eq} in
\lean{Formalisation.C2Cancellation}.
These results apply Mathlib's fixed point counts and classification
of permutations by cycle type,
\lean{Equiv.Perm.isConj_iff_cycleType_eq}, to involutions.
\end{implementation}

The pair $(|X|,|X^\sigma|)$ determines a finite set with involution up
to equivariant bijection. Both entries add under invariant disjoint
unions. Consequently, if $X\simeq Y$ equivariantly and invariant subsets
$X'\subseteq X$, $Y'\subseteq Y$ are equivariantly bijective, then
$X\setminus X'\simeq Y\setminus Y'$ equivariantly. The argument
constructs a bijection of the complements from their two counts.

For two parts exchanged by the involution, a different formula is useful.
Choose a bijection $f:X_0\to Y_0$, and define the map on $X_1$ by
$\tau f\sigma$. It takes $X_1$ to $Y_1$, is inverse to the analogous
map constructed from $f^{-1}$ and commutes with the involutions.
These constructions are in \module{Formalisation.BlockCancellation}
and \module{Formalisation.C2Cancellation}.

\begin{proposition}[Cancellation over a block orbit]
\label{lib:actions:block-cancellation}
Let $I$ be a finite set with an involution, and let $X\to I$ and
$Y\to I$ be equivariant maps of finite sets with involutions.
Write $I_0,\ldots,I_r$ for the orbits in $I$, and let $U_j$ and
$V_j$ be their inverse images in $X$ and $Y$. Suppose that
$X\simeq Y$ equivariantly and that $U_j\simeq V_j$ equivariantly
for $1\leq j\leq r$. Then there is an equivariant bijection
$U_0\to V_0$ preserving the map to $I_0$.
\end{proposition}
\begin{proof}
Cancel the known parts in the total size and fixed point count.
Proposition~\ref{lib:actions:involutions} gives an equivariant
bijection between the remaining parts. If $I_0$ is a singleton,
it already preserves the map to $I_0$.

Otherwise write $I_0=\{i,\sigma i\}$. The involutions exchange
the two fibres on each side, so each fibre has half the size of
its corresponding remaining part. These sizes agree. Choose a
bijection $f:X_i\to Y_i$ and define the map on the other fibre by
$\tau f\sigma$, where $\sigma$ and $\tau$ denote the source and
target involutions. The two formulas give an equivariant bijection
preserving both fibres.
\end{proof}

For an application to blocks, the maps to $I$ assign a block to each
character or weight class. The hypotheses then require equivariant
comparisons on each of the other block orbits.

\begin{implementation}
In \module{Formalisation.BlockCancellation},
\lean{cancel_equivariant_equiv_of_parts} handles the finite
disjoint union of known parts, and
\lean{cancel_exchanged_blocks_preserving} constructs the map on
the remaining pair of exchanged fibres. The resulting map need not
be the restriction of the initially given equivalence $X\simeq Y$.
\end{implementation}

\begin{example}
A set with eight elements and two fixed points has three orbits of size
two. If an invariant part consists of one such orbit, the complement has
the pair of counts $(6,2)$. Any other complement with these counts is
isomorphic as a set with involution.
\end{example}

\section{Fixed representatives in a product}

Let $D\rtimes R$ act on a set $X$. Suppose that
\[
 (D\rtimes R)_x=D_x\rtimes R_x
\]
for some $x\in X$. If $r\in R$ satisfies $rx\in Dx$, then $rx=x$.
Indeed, write $rx=dx$. The element $(d^{-1},r)$ fixes $x$, so the
stabiliser equality implies $r\in R_x$.

\begin{proposition}[A fixed tuple in a product orbit]
Let $n>0$, let $O_0,\ldots,O_{n-1}$ be $D$-orbits in $X$, and let
$T:X\to X$. Suppose a cyclic group $E=\langle\tau\rangle$ acts on
$X^n$ by
\[
 \tau(y_0,\ldots,y_{n-1})=(Ty_{n-1},Ty_0,\ldots,Ty_{n-2}).
\]
Assume that $\tau$ preserves $\prod_jO_j$. Choose $x\in O_0$ such
that
\[
 T^n x=rx\quad\text{for some }r\in R,
 \qquad (D\rtimes R)_x=D_x\rtimes R_x.
\]
Then $(x,Tx,\ldots,T^{n-1}x)$ belongs to $\prod_jO_j$ and is fixed
by $E$.
\end{proposition}
\begin{proof}
Varying the $j$th coordinate of a tuple in the product shows that
$T(O_j)\subseteq O_{j+1}$, with indices taken modulo $n$.
Consequently $T^n x\in O_0=Dx$. The preceding observation gives
$rx=x$, and hence $T^n x=x$.
Successive transport places $T^jx$ in $O_j$.
The formula defining $\tau$ now shows that it fixes the tuple,
so its cyclic group $E$ fixes the tuple as well.
\end{proof}

Apply this construction on each cycle when there are several cycles
of coordinates. The groups $D$ and $R$ may vary from cycle to cycle.
If the actions extend compatibly to $D_{\mathrm{prod}}\rtimes E$,
where $D_{\mathrm{prod}}$ is the product of the coordinate groups,
the resulting tuple $\theta$ satisfies
\[
 (D_{\mathrm{prod}}\rtimes E)_\theta
   =(D_{\mathrm{prod}})_\theta\rtimes E.
\]
Indeed, every element of $E$ fixes $\theta$, and $(d,e)$ fixes
$\theta$ exactly when $d$ does.

\begin{implementation}
\module{ModularRep.ComponentReturnFull} allows separate groups and
orbit data on every cycle.
The theorem \lean{component_return_full_of_product_orbit_stable} in
\lean{ModularRep.ManuscriptVerification.ComponentReturnFull}
derives the conditions on coordinate orbits from stability of the
product orbit and constructs the fixed tuple.
In a representation theoretic application, the product character
sets, their diagonal actions and the action after a full cycle must
satisfy the stated action formulas.
\end{implementation}
